\scanpage{6}
% L1–5
Nous avons déjà \deleted{\unsure{sur}} \inserted{dans la dém de} prop. 1 vérifié que les flèches verticales sont des inclusions de sous-topos [inclusion d'un sous-topos \unsure{ensembliste|ensemblistique|intersection} dans un des facteurs] donc pour prouver que $\Pi(\mathrm{id}_{\widetilde C},v)_*$ est le \inserted{foncteur} image directe pour un morphisme de topos, il suffit de vérifier qu'il en est ainsi de $\Pi(\mathrm{id}_{\widehat C},F')_*$. En effet,

% L6
On aura alors un morphisme de topos

% L7
\diagram{\begin{tikzcd}[ampersand replacement=\&,row sep=small,column sep=small]
\pi(\widetilde C,F)\arrow[rr]\arrow[dr] \& \& \pi(\widehat C,F')\\
\& \pi(\widehat C,F)\arrow[ur] \&
\end{tikzcd}}

% L7–11
, et on voit (en vérifiant sur les foncteurs images directs) qu'il se factorise par le sous topos $\pi(\unsure{\overline C|\widetilde C},F')$, le foncteur image directe associé au \unsure{morphisme} \illegible{1} topos obtenu $\pi(\widetilde C,F)\longrightarrow\pi(\widetilde C,F')$ étant $\pi(\mathrm{id}_{\widetilde C},v)_*$.

% L12
Or $\pi(\mathrm{id}_{\widehat C},\unsure{v)_*}$ \unsure{peut} s'identifier au foncteur

% L13
\[
\underline{\mathrm{Hom}}\,(C,F)\unsure{\longrightarrow|-}\underline{\mathrm{Hom}}\,(C,F')
\]

% L14–19
défini par \deleted{\unsure{$\varphi\mapsto\varphi\circ v^*$}, dont l'adjoint est $\varphi'\mapsto\unsure{\varphi_*}v_*\circ\varphi'$,} $\varphi\longmapsto v_*\circ\varphi$ [car le bifoncteur associé à $\varphi$ est \deleted{\unsure{$\varphi$}} $(X,Y)\to\mathrm{Hom}_F(Y,\varphi(X))$, qu'on \unsure{remplace|compose|transforme} par \deleted{\unsure{transport|Hom}} $(X,Y')\unsure{\longrightarrow|-}\mathrm{Hom}_F(v^*Y',\varphi(X))\simeq\mathrm{Hom}_{F'}(Y',(v_*\varphi)(X))$, et posant $\varphi'=v_*\varphi\colon C\longrightarrow F'$, c'est bien le bifoncteur sur $C\times F'$ associé à $\varphi'$], dont l'adjoint à droite est $\varphi\longmapsto v^*\circ\varphi$, or celui-ci est bien exact à g. \qquad cqfd.

% L20–21
Prenant $v\colon F\longrightarrow P$ (topos ponctuel $\simeq\widehat e$, $e$ catégorie ponctuelle) et notant que $\pi(E,P)\cong E$, on trouve un morphisme can.

% L22
\[
\pi(E,F)\longrightarrow E
\]

% L23
et on trouve de $=$ un morphisme can. de topos

% L24
\[
\pi(E,F)\longrightarrow F
\]

% L25
d'où un morphisme de topos

% L26
\[
\pi(E,F)\longrightarrow E\mathbin{\underset{\mathrm{top}}{\overset{2}{\times}}}F
\]

% L27–28
évidemment 2-fonctoriel en $E,F$. On veut prouver que c'est une \deleted{\unsure{trans}} équivalence, i.e. que $\pi(E,F)$ ``est'' un 2-produit des

% foot
% archive sheet number: 5
